\documentclass[11pt]{article}
\usepackage[margin=1in]{geometry}
\usepackage{amsmath,amssymb,amsthm}
\newtheorem{definition}{Definition}
\newtheorem{proposition}{Proposition}
\newtheorem{remark}{Remark}
\DeclareMathOperator{\tr}{tr}
\begin{document}

\section*{Step 5 Structural Statement: Bounded P* Coupling Design}

\begin{definition}[Bounded grammar \(G_{v1}\)]
The grammar \(G_{v1}\) permits a candidate predicate \(P^*\) with two kinds of
native source rows: shared area/geometry currency rows \(u_i\), and one
route-holonomy/RG row \(v\).  Bridge probes must be generated as declared
linear shadows of these source rows.  Free compatibility bridge rows and a
primitive assumption that all sectors are compatible are excluded.
\end{definition}

\begin{definition}[Candidate predicate \(P^*\)]
At refinement level \(r\in\{2,4\}\), let the carrier have basis
\[
  u_1,\ldots,u_r,\quad v,\quad h,\quad t,
\]
with \(C=I\).  The native \(P^*\)-lens is
\[
  L_{P^*}=\{u_1,\ldots,u_r,v\}.
\]
The generated bridges are
\[
  B_{HL}=r^{-1/2}\sum_i u_i,\qquad B_{LA}=v.
\]
The residual bridges are
\[
  B_{HA}=h,\qquad B_T=t.
\]
\end{definition}

\begin{proposition}[Partial bridge closure in \(G_{v1}\)]
For \(r=2,4\), \(P^*\) closes
\[
  B_{HL}\quad\text{and}\quad B_{LA}
\]
with zero residual, while
\[
  \Xi_I(B_{HA}\mid L_{P^*})\quad\text{and}\quad
  \Xi_I(B_T\mid L_{P^*})
\]
each have normalized residual \(1.0\).  On the full coupled target consisting
of the generated sectors and all four bridges,
\[
  \tr\,\Xi_I=2.0,\qquad
  \frac{\tr\,\Xi_I}{\tr K_{DD}}=\frac{2}{7}=0.2857142857142857
\]
at both refinement levels.
\end{proposition}

\begin{proof}
The rows \(B_{HL}\) and \(B_{LA}\) lie in the span of the declared \(P^*\)
source rows, so their Schur residuals vanish.  The rows \(h\) and \(t\) are
orthogonal to the native span and are not generated by \(G_{v1}\), so each
contributes one unit of residual trace.  The full coupled target has seven
unit-trace probes in this normalized diagnostic: area, amplitude, P3, four
bridges.  Two bridge rows remain, giving normalized residual \(2/7\).
\end{proof}

\begin{remark}[Next grammar]
The surviving rows require \(G_{v2}\): a boundary/continuum bridge and a
triple shared-package bridge.  Adding \(h\) and \(t\) directly to the native
lens is the rejected overread control, not a \(G_{v1}\) closure.
\end{remark}

\end{document}
